\section{Results}
\label{sec:results}

We first compare standard shared-bias ParT with and without edge-conditioned
value messages, then examine the bias scheme and class-wise background
rejection. The main comparisons retain standard ParT pair features.
All results use the same official 20M-jet test sample and are averaged
over three independently trained models per configuration. Reported spreads
are sample standard deviations across training seeds, not confidence intervals
or ensemble results.

\subsection{EdgeValue with standard shared bias}
\label{sec:results-simple-ev}

Table~\ref{tab:edgevalue-baseline} compares the baseline with its direct
EdgeValue extension, retaining the standard four pair features and shared
attention bias. Accuracy increases from $82.819\pm0.100\%$ to
$83.474\pm0.047\%$, a gain of $0.655$ percentage points (pp).
Macro AUROC increases from $98.229\%$ to $98.342\%$, a gain of $0.113$ pp.
The accuracy improvement is positive for each matched seed, ranging from
$0.597$ to $0.726$ pp.

% Generated by tools/generate_tables.py; do not edit by hand.
\begin{table}[htbp]
\centering
\small
\caption{Adding EdgeValue to standard ParT with shared pair bias. Both models use the standard pair features and are trained on 1M jets; evaluation uses the same official 20M-jet test sample. Values are means $\pm$ sample standard deviations across three training seeds. $\Delta$ denotes the difference from ParT in percentage points.}
\label{tab:edgevalue-baseline}
\begin{tabular}{lrrrr}
\toprule
Model & Accuracy (\%) & $\Delta$ (pp) & Macro AUROC (\%) & $\Delta$ (pp) \\
\midrule
ParT & $82.819 \pm 0.100$ & $0.000$ & $98.229 \pm 0.018$ & $0.000$ \\
ParT-EV & $83.474 \pm 0.047$ & $+0.655$ & $98.342 \pm 0.013$ & $+0.113$ \\
\bottomrule
\end{tabular}
\end{table}


This comparison demonstrates that additional pair descriptors and layerwise
bias are not required to obtain an improvement in this setting. The existing
kinematic pair representation becomes more useful when it can contribute to
the vector content of messages as well as their attention weights. The
comparison holds the inputs and bias scheme fixed, but adds the EdgeValue
projections described in Section~\ref{sec:edgevalue}; it is not a
parameter-matched test against a larger bias-only model.

\subsection{Shared versus layerwise bias}
\label{sec:results-bias}

Table~\ref{tab:factorial} compares the four standard-feature configurations
under shared and layerwise bias. Each row compares otherwise matched models
with and without EdgeValue; comparisons between rows measure differences
associated with the bias scheme and its implementation.

% Generated by tools/generate_tables.py; do not edit by hand.
\begin{table}[t]
\centering
\small
\caption{EdgeValue with shared and layerwise pair bias on the official 20M-jet test sample after training on 1M jets. All four configurations use the standard ParT pair features. Within each row, only EdgeValue is switched on or off. Values are means $\pm$ sample standard deviations over three seeds; EV gain is the difference of the two means in percentage points.}
\label{tab:factorial}
\begin{tabular*}{\textwidth}{@{\extracolsep{\fill}}lrrr@{}}
\multicolumn{4}{@{}l}{\textbf{(a) Accuracy (\%)}} \\
\toprule
Pair bias & Without EV & With EV & EV gain (pp) \\
\midrule
Shared & $82.819 \pm 0.100$ & $83.474 \pm 0.047$ & $+0.655$ \\
Layerwise & $82.859 \pm 0.055$ & $83.438 \pm 0.052$ & $+0.579$ \\
\bottomrule
\end{tabular*}
\par\vspace{0.8em}
\begin{tabular*}{\textwidth}{@{\extracolsep{\fill}}lrrr@{}}
\multicolumn{4}{@{}l}{\textbf{(b) Macro AUROC (\%)}} \\
\toprule
Pair bias & Without EV & With EV & EV gain (pp) \\
\midrule
Shared & $98.229 \pm 0.018$ & $98.342 \pm 0.013$ & $+0.113$ \\
Layerwise & $98.233 \pm 0.008$ & $98.330 \pm 0.006$ & $+0.096$ \\
\bottomrule
\end{tabular*}
\end{table}


With standard pair features, moving from shared to layerwise bias changes the
mean accuracy by only $+0.040$ pp without EdgeValue. With EdgeValue, the
layerwise configuration reaches $83.438\%$, slightly below the shared-bias
model's $83.474\%$ ($-0.036$ pp). Neither ordering holds for all three seeds:
layerwise bias improves accuracy in two seeds without EdgeValue and in one
seed with EdgeValue. These small mean differences do not show a consistent
advantage for the layerwise configuration with standard inputs.

By contrast, adding EdgeValue improves accuracy under both bias schemes:
$+0.655$ pp with shared bias and $+0.579$ pp with layerwise bias. Both
comparisons improve in all three seeds. The benefit of edge-conditioned
messages therefore persists across these two implementations rather than
depending on the layerwise bias choice. As discussed in
Section~\ref{sec:training}, shared-versus-layerwise comparisons also change
trimming and pair-normalization handling, so they should not be interpreted
as isolating bias sharing alone.

\subsection{Class-wise background rejection}
\label{sec:results-rejection}

Table~\ref{tab:rejection-main} compares the four standard-feature
configurations at the JetClass benchmark operating points. The columns
pair each ParT-EV variant with its corresponding bias-only control,
retaining the same pair inputs and bias scheme within each comparison.
Complete numerical results are provided in
Appendix~\ref{sec:complete-results}.

% Generated by tools/generate_tables.py; do not edit by hand.
\begin{table}[t]
\centering
\footnotesize
\setlength{\tabcolsep}{4pt}
\renewcommand{\arraystretch}{1.15}
\caption{QCD background rejection on the official 20M-jet test sample at the JetClass benchmark signal efficiencies~\cite{Qu2022ParticleTransformer}. The first two model columns compare standard shared-bias ParT without and with EV; the last two compare layerwise models without and with EV. All models use standard ParT pair features and are trained on 1M jets. Values are means $\pm$ sample standard deviations over three seeds; higher is better. }
\label{tab:rejection-main}
\begin{tabular*}{\textwidth}{@{\extracolsep{\fill}}lc*{4}{c}@{}}
\toprule
 & & \multicolumn{2}{c}{Shared bias} & \multicolumn{2}{c}{Layerwise bias} \\
\cmidrule(lr){3-4}\cmidrule(lr){5-6}
Signal & $\epsilon_s$ (\%) & ParT & ParT-EV & ParT & ParT-EV \\
\midrule
$H\!\to b\bar b$ & 50 & $5271 \pm 242$ & $5674 \pm 152$ & $5191 \pm 195$ & $5408 \pm 248$ \\
$H\!\to c\bar c$ & 50 & $1741 \pm 31$ & $2007 \pm 15$ & $1792 \pm 84$ & $2020 \pm 143$ \\
$H\!\to gg$ & 50 & $90.3 \pm 0.8$ & $94.1 \pm 0.8$ & $91.2 \pm 0.7$ & $93.7 \pm 1.7$ \\
$H\!\to4q$ & 50 & $665.0 \pm 22.3$ & $813.5 \pm 10.0$ & $683.6 \pm 32.5$ & $788.6 \pm 24.0$ \\
$H\!\to q\bar q'\ell\nu$ & 99 & $1564 \pm 23$ & $1832 \pm 32$ & $1535 \pm 126$ & $1855 \pm 106$ \\
$Z\!\to q\bar q$ & 50 & $229.8 \pm 3.2$ & $249.5 \pm 4.9$ & $234.8 \pm 5.0$ & $249.1 \pm 4.4$ \\
$W\!\to q\bar q'$ & 50 & $284.7 \pm 7.8$ & $315.2 \pm 9.4$ & $292.1 \pm 2.7$ & $311.7 \pm 2.3$ \\
$t\!\to b q\bar q'$ & 50 & $5833 \pm 464$ & $7395 \pm 263$ & $5629 \pm 40$ & $6932 \pm 761$ \\
$t\!\to b\ell\nu$ & 99.5 & $3903 \pm 161$ & $4207 \pm 182$ & $3541 \pm 348$ & $4267 \pm 378$ \\
\bottomrule
\end{tabular*}
\end{table}


The standard shared-bias EdgeValue extension increases the mean background
rejection for all nine signal classes at these operating points. Among the
hadronic classes at 50\% signal efficiency, the relative increases range
from $4.2\%$ for $H\to gg$ to $26.8\%$ for fully hadronic top jets.
For $H\to4q$, rejection rises from $665.0$ to $813.5$ ($+22.3\%$),
while for $H\to b\bar b$ it rises from $5271$ to $5674$ ($+7.6\%$).
Thus, the improvement observed in multiclass accuracy is accompanied by
better signal-versus-QCD discrimination across the benchmark categories.

With layerwise bias, adding EdgeValue likewise improves mean benchmark
rejection in all nine classes relative to its matching bias-only control.
The gains include $15.4\%$ for $H\to4q$ and $23.2\%$ for fully
hadronic top jets at 50\% efficiency. Thus, the class-wise improvements
also persist under both bias schemes, although their sizes differ by class.

At the stricter 30\% operating point
(Table~\ref{tab:rejection-30}), the standard shared-bias EdgeValue
model improves mean rejection over baseline in all seven hadronic classes,
including $28.1\%$ for $H\to4q$ and $35.6\%$ for fully hadronic top jets.
These comparisons show that the benefits extend to a lower-efficiency regime.
All rejection gains quoted here are relative changes in mean rejection,
distinct from the percentage-point changes used for accuracy and macro
AUROC; point-estimate differences alone do not establish statistical
significance.

\paragraph{Beyond standard pair features.}
Appendix~\ref{sec:richer-relations} examines richer physics-inspired pair
representations and finds positive EdgeValue accuracy gains under both bias
schemes. These supplementary results motivate a broader study of which
relations best complement edge-conditioned value messages, without requiring
those additional descriptors for the gains demonstrated in the main study.
